\section{Introduction}
Human activity recognition (HAR) aims to infer behaviors such as walking, sitting, and standing from wearable sensor streams. While the UCI HAR benchmark provides labels for evaluation, we treat the learning problem as fully unsupervised: labels are used only after training to measure how well the discovered latent states align with real activities.

Our project compares two inference strategies applied to the same underlying Gaussian hidden Markov model. The frequentist baseline estimates parameters with EM, while the Bayesian model places priors over the same transition and emission parameters and samples from the posterior with Gibbs updates. This distinction clarifies the core scientific question from the proposal: does Bayesian inference improve the practical behavior of the same sequential model, rather than changing the model class itself?

The project also addresses the concern raised in the proposal feedback about robustness. Here, robustness is defined concretely as the model's ability to maintain clustering quality when Gaussian noise is injected into the held-out test features, together with held-out log-likelihood as a measure of probabilistic fit. Uncertainty quantification matters in this setting because Bayesian inference yields posterior distributions over transitions and emissions, making it possible to inspect whether apparently persistent activities are estimated with confidence or with substantial variability.
